\documentclass[a4paper,12pt]{article}
\usepackage[utf8]{inputenc}
\usepackage[english, russian]{babel}
\usepackage{amsmath, amssymb, amsthm}
\usepackage{enumitem}%оформление списков
\usepackage{varwidth}
\usepackage{graphicx}
\usepackage{float}
\usepackage{indentfirst}
\usepackage{url}
\usepackage{seqsplit}
\usepackage{pdflscape}
\usepackage{tabularx}
\usepackage{xcolor}
\usepackage{listings}
\usepackage{mathtools}
\usepackage{hyperref}
\usepackage[backend=biber, style=numeric, sorting=none]{biblatex}
\usepackage{titlesec}
\usepackage{appendix}
\usepackage{alltt}
\usepackage[hidelinks]{hyperref}
\usepackage{booktabs}

\addbibresource{my_libs.bib}


\usepackage[a4paper, top=20mm, left=30mm, right=20mm, bottom=20mm]{geometry}
\newcommand{\dim}{\mathrm{dim}}
\newcommand{\lin}{\mathrm{lin}}
\newcommand{\diag}{\mathrm{diag}}
\newcommand{\R}{\mathbb{R}}
\newcommand{\Rn}{\mathbb{R^n}}

\lstset{
    language=C++,
    basicstyle=\ttfamily\small,
    keywordstyle=\color{blue},
    commentstyle=\color{green!60!black},
    stringstyle=\color{red},
    numbers=left,
    numberstyle=\tiny\color{gray},
    stepnumber=1,
    breaklines=true,
    breakatwhitespace=true,
    tabsize=4,
    showspaces=false,
    showstringspaces=false
}

\newtheoremstyle{indented}% отступ
{3pt}% над
{3pt}% под
{\itshape}% шрифт тела, переключатель
{\parindent}% отступ
{\bfseries}% шрифт заголовка
{.}% точка после заголовка
{.5em}% отступ от заголовка
{}% head spec (empty = ‘normal’)

\theoremstyle{indented}
\newtheorem{theorem}{Теорема}
\newtheorem*{theorem*}{Теорема}
\newtheorem{definition}{Определение}
\newtheorem*{definition*}{Определение}
\newtheorem{following}{Следствие}
\newtheorem*{following*}{Следствие}
\newtheorem{statement}{Утверждение}
\newtheorem*{statement*}{Утверждение}

\renewcommand{\thesection}{\arabic{section}.}
\renewcommand{\thesubsection}{\arabic{section}.\arabic{subsection}.}
\renewcommand{\thesubsubsection}{\arabic{section}.\arabic{subsection}.\arabic{subsubsection}.}

\setcounter{section}{0}
\setcounter{subsection}{0}
\setcounter{equation}{0}

\newcommand{\setupappendix}{
  \setcounter{section}{0} % Сбрасываем счетчик разделов для приложений
  \renewcommand{\thesection}{\Alph{section}} % Буквенная нумерация (А, Б, В...)

  % Настройка отображения самого заголовка в тексте (по правому краю)
  \titleformat{\section}[display]
    {\normalfont\Large\bfseries\filleft}
    {Приложение \thesection}
    {10pt}
    {\Large}
}

\begin{document}
\pagenumbering{arabic}
\setcounter{page}{1}
\pagestyle{empty}
\begin{center}
\textbf{Министерство науки и высшего образования Российской Федерации}
\\
\vspace{0.5cm}
\textbf{Федеральное государственное бюджетное образовательное\\
учреждение высшего образования} \\
\textbf{«Ярославский государственный университет им. П.Г. Демидова»}

\vspace{0.5cm}

Кафедра математического анализа
\vfill

\begin{flushright}
  Сдано на кафедру\\
  « \underline{\hspace{0.8cm}} » \underline{\hspace{3cm}} 2026 г.\\
  Заведующий кафедрой \\
  \underline{д. ф.-м. н., доцент }\\
  \underline{\hspace{3.4cm}} Невский М.В.\\
\end{flushright}
\vfill

Выпускная квалификационная работа магистра \\

\vspace{0.5cm}

\textbf{Задача о минимальной норме интерполяционного проектора} \\
(Направление подготовки магистров 01.03.02 Прикладная математика и информатика)
\vfill


\bigskip


\end{center}

\begin{flushright}
Научный руководитель\\
\underline{к. ф.-м. н., доцент }\\
\underline{\hspace{3.4cm}} Ухалов А.Ю.\\
« \underline{\hspace{0.8cm}} » \underline{\hspace{3cm}} 2026 г.\\
\vspace{1.5cm}
Студент группы \underline{ПМИ-21МО}\\
\underline{\hspace{3.5cm}} Зайков И.В. \\
« \underline{\hspace{0.8cm}} » \underline{\hspace{3cm}} 2026 г.\\
\end{flushright}

\vspace{0.5cm}
\begin{center}
Ярославль 2026 г.
\end{center}
\newpage

\section*{Реферат}
\addcontentsline{section}{Реферат}

Объем 51 страница, 5 параграфов, 2 таблицы, 3 приложения.
\par
Изучается интерполяция непрерывной функции $f(x)$, определенной на кубе
$Q_n = [0, 1]^n$ , с помощью полиномов от $n$ переменных степени не выше 1. Ставится
задача оценивания минимальной нормы проектора для $n = 31$. Верхняя оценка $\theta_{31}$
получается из рассмотрения проектора с узлами интерполяции, совпадающими с
вершинами симплекса максимального объема в $Q_n$, построенного на основе матрицы Адамара. Получена новая оценка $\theta_{31} \leq 5.0$.
\par
\textbf{Ключевые слова: интерполяция, интерполяционный проектор, норма, матрица Адамара, симплекс максимального объема.}

\newpage

\pagestyle{plain}
\tableofcontents

\newpage

% ============================================================================
\section*{Введение}
\addcontentsline{toc}{section}{Введение}

В выпускной квалификационной работе рассматривается задача о минимальной норме интерполяционного проектора. Изучается интерполяция непрерывной функции $f(x)$, определённой на кубе $Q_n = [0,1]^n$, с помощью полиномов от $n$ переменных степени не выше $1$.
\par
Для построения интерполяционного полинома необходимо задать $n+1$ узел — вершины невырожденного симплекса $S \subset Q_n$.
\par
Интерполяционный проектор $P$ ставит в соответствие функции $f$ полином $P f$, совпадающий с $f$ в узлах. Норма оператора $P$ характеризует устойчивость интерполяции к погрешностям входных данных: чем меньше норма, тем ближе интерполяционный полином к наилучшему приближению.

Минимальная норма интерполяционного проектора $\theta_n$ определяется как наименьшее возможное значение $\|P\|$ при условии, что узлы интерполяции принадлежат кубу $Q_n$ и образуют невырожденный симплекс. Задача нахождения $\theta_n$ является классической в теории приближений и имеет важное практическое значение, поскольку позволяет выбрать расположение узлов, минимизирующее множитель $(1 + \|P\|)$ в неравенстве Лебега и, следовательно, дающее наилучшую гарантированную оценку погрешности интерполяции.

В настоящей работе для получения верхней оценки $\theta_n$ используется подход, основанный на рассмотрении симплексов максимального объёма, вписанных в гиперкуб $Q_n$. Такие симплексы могут быть построены из матриц $(-1/1)$ порядка $n+1$, определители которых максимальны, то есть из матриц Адамара. Норма проектора для симплекса вычисляется по формуле
\[
\|P\| = \max_{x \in \operatorname{ver}(Q_n)} \sum_{j=1}^{n+1} |\lambda_j(x)|,
\]
где $\lambda_j$ — базисные многочлены Лагранжа (барицентрические координаты). Для нахождения максимума требуется перебрать все $2^n$ вершин гиперкуба.
\par
Для размерности $n = 31$ было исследовано шесть типов матриц Адамара порядка $32$ (\texttt{syl}, \texttt{pal}, \texttt{t1}, \texttt{t2}, \texttt{t3}, \texttt{t4}). Для каждого типа построен симплекс максимального объёма $S \subset Q_{31}$ и вычислена норма интерполяционного проектора $\|P_S\|$. В качестве оценки сверху для $\theta_{31}$ принимается минимум норм по всем рассмотренным симплексам:
\[
\theta_{31} \le \min_{S} \|P_S\|.
\]
\par
Была разработана программа на языке C++ с использованием библиотеки Eigen для линейно-алгебраических операций. Для ускорения вычислений использовались параллельные вычисления на $16$ потоках, что позволило завершить полный перебор за $290$ секунд (около $5$ минут).
\par
Для верификации программной реализации был проведён ряд проверок:
\begin{itemize}
    \item проверка интерполяционных свойств $\lambda_j(x^{(k)}) = \delta_{jk}$ с точностью $10^{-10}$;
    \item верификация на пространствах $\mathbb{R}^n$ меньших размерностей путём сравнения с известными из литературы результатами ($n = 0, 1, 3, 7, 11, 15, 19, 23, 27$);
    \item повторная проверка найденного максимума после завершения перебора.
\end{itemize}

В результате работы получена верхняя оценка минимальной нормы интерполяционного проектора для размерности $n = 31$:
\[
\theta_{31} \le 5.000.
\]

\newpage

% ============================================================================
\section{Используемые обозначения}

Везде далее предполагается, что $n$ — натуральное число ($n \in \mathbb{N}$). Обозначим через $Q_n$ единичный куб в $\mathbb{R}^n$:
\[
Q_n = [0,1]^n = \{(x_1, \dots, x_n) : 0 \le x_i \le 1,\ i = 1,\dots,n\}.
\]

Через $C(Q_n)$ обозначим пространство непрерывных функций $f : Q_n \to \mathbb{R}$ с равномерной нормой
\[
\|f\|_{C(Q_n)} = \max_{x \in Q_n} |f(x)|.
\]

Под $\Pi_1(\mathbb{R}^n)$ понимается совокупность многочленов от $n$ переменных степени $\le 1$, т.е. линейная оболочка
\[
\Pi_1(\mathbb{R}^n) = \operatorname{span}\{1, x_1, x_2, \dots, x_n\}.
\]

Симплекс $S \subset \mathbb{R}^n$ с вершинами $x^{(1)}, \dots, x^{(n+1)}$ обозначается
\[
S = \operatorname{conv}\{x^{(1)}, \dots, x^{(n+1)}\}.
\]

% ============================================================================
\section{Интерполяционный проектор}

\subsection{Постановка задачи интерполяции}

В \cite{nevsky_geo} описывается задача линейной интерполяции на $n$-мерном кубе.
\par
Пусть $Q_n = [0,1]^n$ — единичный куб в $\mathbb{R}^n$, а $C(Q_n)$ — пространство непрерывных функций $f: Q_n \to \mathbb{R}$ с равномерной нормой
\[
\|f\|_{C(Q_n)} = \max_{x \in Q_n} |f(x)|.
\]

Рассмотрим интерполяцию функций из $C(Q_n)$ полиномами от $n$ переменных степени не выше $1$. Пространство таких полиномов обозначим через $\Pi_1(\mathbb{R}^n)$:
\[
\Pi_1(\mathbb{R}^n) = \operatorname{span}\{1, x_1, x_2, \dots, x_n\}.
\]

Для построения интерполяционного полинома необходимо задать $n+1$ узел
\[
x^{(1)}, x^{(2)}, \dots, x^{(n+1)} \in Q_n,
\]
образующих невырожденный симплекс $S$. Требуется найти полином $p \in \Pi_1(\mathbb{R}^n)$ такой, что
\[
p(x^{(j)}) = f(x^{(j)}), \quad j = 1, \dots, n+1.
\]

\subsection{Базисные многочлены Лагранжа и барицентрические координаты}

В \cite{nevsky_geo} описан способ построения базисных многочленов Лагранжа с помощью барицентрических координат, определённых относительно вершин симплекса.

Обозначим вершины симплекса $S$ через
\[
x^{(j)} = \left(x_1^{(j)}, x_2^{(j)}, \dots, x_n^{(j)}\right), \quad 1 \le j \le n+1.
\]

Рассмотрим матрицу $A$ размера $(n+1) \times (n+1)$, в которой \textbf{строки} соответствуют вершинам симплекса, дополненным единицей в последнем столбце:
\[
A = \begin{pmatrix}
x_1^{(1)} & x_2^{(1)} & \dots & x_n^{(1)} & 1 \\
x_1^{(2)} & x_2^{(2)} & \dots & x_n^{(2)} & 1 \\
\vdots & \vdots & & \vdots & \vdots \\
x_1^{(n+1)} & x_2^{(n+1)} & \dots & x_n^{(n+1)} & 1
\end{pmatrix}.
\]

Столбец единиц введён для того, чтобы учесть условие нормировки барицентрических координат $\sum_{j=1}^{n+1} \lambda_j(x) = 1$. Переход к однородным координатам позволяет записать систему для нахождения барицентрических координат в матричной форме.
\par
Рассмотрим матрицу $A$ размера $(n+1) \times (n+1)$, строки которой соответствуют вершинам симплекса, дополненным единицей в последнем столбце:
\[
A = \begin{pmatrix}
x_1^{(1)} & x_2^{(1)} & \dots & x_n^{(1)} & 1 \\
x_1^{(2)} & x_2^{(2)} & \dots & x_n^{(2)} & 1 \\
\vdots & \vdots & & \vdots & \vdots \\
x_1^{(n+1)} & x_2^{(n+1)} & \dots & x_n^{(n+1)} & 1
\end{pmatrix}.
\]
\par
Тогда для точки $x = (x_1, \dots, x_n)$ возможно найти ее барицентрические координаты $\lambda(x) = (\lambda_1(x), \dots, \lambda_{n+1}(x))^T$ из системы
\[
A \lambda(x) = b(x),
\]
где $b(x) = (x_1, \dots, x_n, 1)^T$ — расширенный вектор точки. Поскольку симплекс невырожден, матрица $A$ обратима, и решение системы имеет вид
\[
\lambda(x) = A^{-1} b(x).
\]
\par
Пусть $\Delta$ - определитель матрицы $A$. Рассмотрим $\Delta_j(x)$ определитель, который получается из $\Delta$ заменой $j$-й строки на строку $(x_1, x_2, \dots, x_n, 1)$. Тогда функции
\[
\lambda_j(x) = \frac{\Delta_j(x)}{\Delta}, \quad j = 1, \dots, n+1,
\]
являются также барицентрическими координатами точки $x$ относительно симплекса $S$. Они обладают следующими свойствами:
\begin{enumerate}
    \item $\lambda_j(x) \in \Pi_1(\mathbb{R}^n)$, то есть являются линейными функциями;
    \item $\lambda_j(x^{(k)}) = \delta_{jk}$, где $\delta_{jk}$ — символ Кронекера;
    \item $\sum_{j=1}^{n+1} \lambda_j(x) = 1$ для всех $x \in \mathbb{R}^n$.
\end{enumerate}

Таким образом, барицентрические координаты $\lambda_j(x)$ являются \textbf{базисными многочленами Лагранжа}, соответствующими симплексу $S$. Коэффициенты этих многочленов составляют $j$-й столбец обратной матрицы $A^{-1}$. Действительно, пусть $A^{-1} = (l_{ij})_{i,j=1}^{n+1}$. Тогда
\[
\lambda_j(x) = l_{1j} x_1 + l_{2j} x_2 + \dots + l_{nj} x_n + l_{n+1,j}.
\]

Любой полином $p \in \Pi_1(\mathbb{R}^n)$ может быть разложен по этому базису:
\[
p(x) = \sum_{j=1}^{n+1} p(x^{(j)}) \lambda_j(x). \tag{1}
\]

\subsection{Интерполяционный проектор}

Для заданной функции $f \in C(Q_n)$ определим отображение $P: C(Q_n) \to \Pi_1(\mathbb{R}^n)$ формулой
\[
P f(x) = \sum_{j=1}^{n+1} f(x^{(j)}) \lambda_j(x). \tag{2}
\]

Из свойства $\lambda_j(x^{(k)}) = \delta_{jk}$ следует, что
\[
P f(x^{(j)}) = f(x^{(j)}), \quad j = 1, \dots, n+1,
\]
то есть полином $P f$ интерполирует функцию $f$ в узлах $x^{(j)}$. Отображение $P$ называется \textbf{интерполяционным проектором}.

Оператор $P$ линеен и обладает свойством идемпотентности: для любого полинома $p \in \Pi_1(\mathbb{R}^n)$ выполнено $P p = p$.

\subsection{Норма интерполяционного проектора}

Норма оператора $P$ определяется стандартным образом:
\[
\|P\| = \sup_{\|f\|_{C(Q_n)} \le 1} \|P f\|_{C(Q_n)}.
\]

В \cite{nevsky_geo} показано, что для гиперкуба $Q_n$ максимум суммы модулей базисных функций достигается в одной из его вершин, т.е. справедлива формула:
\[
\|P\| = \max_{x \in \operatorname{ver}(Q_n)} \sum_{j=1}^{n+1} |\lambda_j(x)|. \tag{3}
\]

\subsection{Неравенство Лебега}

% todo: ссылку на литературу
Пусть $E_n(f)$ — наилучшее приближение функции $f$ полиномами из $\Pi_1(\mathbb{R}^n)$:
\[
E_n(f) = \min_{p \in \Pi_1(\mathbb{R}^n)} \|f - p\|_{C(Q_n)}.
\]

Для интерполяционного проектора $P$ справедливо \textbf{неравенство Лебега}:
\[
\|f - P f\|_{C(Q_n)} \le (1 + \|P\|) \cdot E_n(f). \tag{4}
\]

Это неравенство показывает, что погрешность интерполяции не превосходит наилучшего приближения, умноженного на множитель $(1 + \|P\|)$. Таким образом, чем меньше норма проектора, тем ближе интерполяционный полином к наилучшему приближению.

\subsection{Минимальная норма проектора}

Обозначим через $\theta_n$ минимальную норму интерполяционного проектора при условии, что узлы интерполяции принадлежат кубу $Q_n$ и образуют невырожденный симплекс:
\[
\theta_n = \min_{\substack{x^{(j)} \in Q_n \\ \det(A) \neq 0}} \|P\|.
\]
\par
Задача о нахождении $\theta_n$ является классической в теории приближений. Она сводится к поиску такого набора узлов (симплекса $S \subset Q_n$), который минимизирует правую часть (3). Отметим, что выражение под знаком минимума зависит от $n(n+1)$ переменных (координат $n+1$ вершины симплекса), а вычисление $\lambda_j(x)$ требует обращения матрицы $A$, что делает задачу крайне трудоёмкой для больших $n$.
\par
Задача настоящей работы состоит в получении верхней оценки для $\theta_{31}$. Для этого рассматривается конечное множество симплексов $\mathcal{S}$, построенных по шести матрицам Адамара порядка $32$. Для каждого $S \in \mathcal{S}$ вычисляется норма $\|P_S\|$, после чего выбирается минимальное значение:
\[
\theta_{31} \le \min_{S \in \mathcal{S}} \|P_S\|
\]

\subsection{Оценки нормы проектора}

Для величины $\theta_n$ известны следующие общие оценки (см. \cite{nevsky_geo}):
\[
\frac{1}{4}\sqrt{n} < \theta_n < 3\sqrt{n}, \qquad 3 - \frac{4}{n+1} \le \theta_n.
\]

Если $S$ — симплекс максимального объёма в $Q_n$, а $P$ — интерполяционный проектор с узлами в вершинах $S$, то справедливо асимптотическое равенство
\[
\|P\| \asymp \theta_n,
\]
где запись $f(n) \asymp g(n)$ означает, что существуют константы $c_1, c_2 > 0$, не зависящие от $n$, такие что
\[
c_1 g(n) \le f(n) \le c_2 g(n) \quad \text{для всех достаточно больших } n.
\]

Таким образом, норма проектора по узлам в вершинах симплекса максимального объёма и величина $\theta_n$ имеют одинаковый порядок роста. Это свойство лежит в основе метода уточнения оценок для $\theta_n$, применяемого в настоящей работе. В качестве верхней границы $\theta_n$ можно взять значение нормы проектора с узлами в вершинах симплекса максимального объёма в $Q_n$.
\par
В \cite{nevsky_ukhalov_sel_problems} приведены оценки $\theta_n$ для различных размерностей. Наиболее точные известные оценки для $1 \le n \le 26$ систематизированы в таблице ~\ref{tab:theta_estimates}.
% ============================================================================
\section{Матрицы Адамара}

\subsection{Определение и основные свойства}

\textbf{($-1/1$)-матрицей} называется матрица, каждый элемент которой равен $-1$ или $1$. \textbf{Матрицей Адамара} порядка $n$ называется квадратная $(-1/1)$-матрица $H$ размера $n \times n$, удовлетворяющая условию:
\[
H H^T = n I_n,
\]
где $I_n$ — единичная матрица. Из этого свойства непосредственно следует, что строки (и столбцы) матрицы Адамара попарно ортогональны. Для определителя матрицы Адамара справедливо соотношение:
\[
|\det H| = n^{n/2}.
\]

Подробная информация о матрицах Адамара содержится в книге \cite{hall_comb_1970}

Матрицы Адамара существуют не для всех порядков. Необходимым условием существования является $n = 1$, $n = 2$ или $n$, кратное $4$. Достаточность этого условия составляет содержание \textbf{гипотезы Адамара}, которая до настоящего времени не доказана. Тем не менее, для огромного количества порядков, кратных $4$, матрицы Адамара построены. В частности, для порядка $32 = 4 \cdot 8$ матрицы Адамара существуют.

\subsection{Количество неэквивалентных матриц Адамара}

Две матрицы Адамара называются \textbf{эквивалентными}, если одна может быть получена из другой с помощью перестановок строк и столбцов и умножения строк и столбцов на $-1$. Количество неэквивалентных матриц Адамара для малых порядков приведено в таблице~\ref{tab:hadamard_counts}.

% todo: сформулировать
Из таблицы ~\ref{tab:hadamard_counts} видно, что количество матриц быстро растёт с увеличением порядка. Для $n = 32$ оно составляет уже более тринадцати миллионов, что делает полный перебор всех симплексов максимального объёма в $\mathbb{R}^{31}$ практически невозможным.

\subsection{Методы построения матриц Адамара}

Существует несколько различных методов построения матриц Адамара. В настоящей работе используются матрицы порядка $32$, взятые из открытой библиотеки \cite{sloane_hadamard}. Перечень использованных типов матриц приведён ниже:

\subsubsection{Метод Сильвестра (тип `syl`)}

Для порядков, являющихся степенью двойки, матрицы Адамара могут быть построены рекурсивным методом Сильвестра. Рекуррентное соотношение имеет вид:
\[
H_{1} = (1), \qquad H_{2n} = \begin{pmatrix}
H_n & H_n \\
H_n & -H_n
\end{pmatrix}.
\]

Например, для $n = 2$ получаем:
\[
H_2 = \begin{pmatrix}
1 & 1 \\
1 & -1
\end{pmatrix},
\]
для $n = 4$:
\[
H_4 = \begin{pmatrix}
1 & 1 & 1 & 1 \\
1 & -1 & 1 & -1 \\
1 & 1 & -1 & -1 \\
1 & -1 & -1 & 1
\end{pmatrix}.
\]

Матрицы, построенные этим методом, обозначаются суффиксом \texttt{syl} (например, \texttt{had.32.syl}).

\subsubsection{Метод Пэли (тип `pal`)}

Конструкция Пэли (Paley construction) позволяет строить матрицы Адамара для порядков $n = q + 1$, где $q$ — степень нечётного простого числа. В основе метода лежат свойства квадратичных вычетов в конечных полях. Для порядка $32 = 31 + 1$ существует матрица Адамара типа Пэли, обозначаемая суффиксом \texttt{pal} (например, \texttt{had.32.pal}).

\subsubsection{Другие типы матриц порядка 32}

Для порядка $32$ кроме матриц типов \texttt{syl} и \texttt{pal} существуют также матрицы, обозначаемые \texttt{t1}, \texttt{t2}, \texttt{t3}, \texttt{t4} (например, \texttt{had.32.t1}, \texttt{had.32.t2}, \texttt{had.32.t3}, \texttt{had.32.t4}). Эти матрицы получены методом тензорного произведения. В настоящей работе исследуются все перечисленные типы матриц для сравнения значений норм интерполяционных проекторов, получаемых на соответствующих симплексах.

\subsection{Переход от $(-1/1)$-матриц к $(0/1)$-матрицам максимального определителя}

Как показано в \cite{nevsky_geo}, симплекс максимального объёма в гиперкубе $Q_n = [0,1]^n$ может быть получен из $(0/1)$-матрицы порядка $n$, имеющей максимальный определитель. В качестве первых $n$ вершин такого симплекса можно взять строки этой матрицы, а в качестве последней $(n+1)$-й вершины добавить нулевую строку.
Максимальный $(0/1)$-определитель порядка $n$ может быть получен из максимального $(-1/1)$-определителя порядка $n+1$ с помощью следующей процедуры. Пусть $A$ — максимальный $(-1/1)$-определитель порядка $n+1$.

\begin{enumerate}
    \item Каждый столбец $A$, начинающийся с $-1$, умножается на $-1$.
    \item Каждая строка новой матрицы, начинающаяся с $-1$, также умножается на $-1$.
\end{enumerate}

Матрицу порядка $n+1$, которая получится в результате выполнения шагов 1–2, обозначим через $B$. У этой $(-1/1)$-матрицы первый столбец и первая строка целиком состоят из $1$.

\begin{enumerate}
    \setcounter{enumi}{2}
    \item Пусть $C$ — подматрица порядка $n$ матрицы $B$, стоящая в строках и столбцах с номерами $2, \dots, n+1$. В этой матрице производится следующая замена элементов: $1$ заменяется на $0$, а $-1$ на $1$. Получившуюся $(0/1)$-матрицу порядка $n$ обозначим $D$.
\end{enumerate}

Полученный таким образом определитель $D$ является максимальным $(0/1)$-определителем порядка $n$.

\subsection{Построение симплекса максимального объёма в $\mathbb{R}^{31}$}

В настоящей работе для размерности $n = 31$ в качестве исходных берутся матрицы Адамара порядка $32$ следующих типов: \texttt{syl}, \texttt{pal}, \texttt{t1}, \texttt{t2}, \texttt{t3}, \texttt{t4}. Для каждой из них выполняется описанная выше процедура перехода к $(0/1)$-матрице $D$ порядка $31$, после чего строки $D$ становятся первыми $31$ вершиной симплекса, а в качестве последней, $32$-й вершины добавляется нулевая строка:
\[
V = \{0, \operatorname{row}_1(D), \operatorname{row}_2(D), \dots, \operatorname{row}_{31}(D)\} \subset [0,1]^{31}.
\]

Таким образом, для каждого типа матрицы Адамара получается соответствующий симплекс максимального объёма, вписанный в гиперкуб $[0,1]^{31}$. Все вершины таких симплексов являются вершинами гиперкуба. Полученные симплексы используются в настоящей работе для вычисления верхних оценок минимальной нормы интерполяционного проектора $\theta_{31}$.

\subsection{Замечание о количестве симплексов}

Для размерности $n = 31$ количество неэквивалентных матриц Адамара порядка $32$ составляет $13\,710\,027$ \cite{kharaghani2013hadamard}. Это означает, что существует такое же количество различных симплексов максимального объёма, вписанных в гиперкуб $Q_{31}$. Значения норм интерполяционных проекторов для разных симплексов могут различаться.
\par
В настоящей работе для исследования используются шесть матриц Адамара порядка $32$, доступных на сайте библиотеки Нила Слоуна \cite{sloane_hadamard}:
\[
\texttt{had.32.syl},\quad \texttt{had.32.pal},\quad \texttt{had.32.t1},\quad \texttt{had.32.t2},\quad \texttt{had.32.t3},\quad \texttt{had.32.t4}.
\]

Полный перебор всех $13\,710\,027$ симплексов является крайне трудоёмкой задачей и не входит в рамки настоящей работы. Исследование ограничивается указанными шестью представителями, что позволяет получить содержательные оценки при разумных вычислительных затратах.
% ============================================================================
\section{Описание программной реализации}

Исходный код программы доступен в открытом репозитории по адресу:
\url{https://gitflic.ru/project/zaykov_ivan/minimal_norm_calculator.git} (или \url{https://github.com/Ivan-Zaykov/minimal_norm_calculator.git}).

Программная реализация выполнена на языке C++ (стандарт C++17). Для линейно-алгебраических операций используется библиотека Eigen \cite{eigen_lib}. Архитектура программы позволяет легко адаптировать её для вычисления нормы интерполяционного проектора для любых симплексов максимального объёма, не только построенных из матриц Адамара. Сборка проекта осуществляется с помощью CMake.

\subsection{Общая архитектура}

Программа построена в соответствии с принципами SOLID. Основные классы:

\begin{itemize}
    \item \texttt{ISimplexProvider} — интерфейс для получения вершин симплекса;
    \item \texttt{SilvesterMethodProvider} — построение симплекса из матрицы Адамара методом Сильвестра;
    \item \texttt{HadamardFileProvider} — построение симплекса по матрице Адамара из файла;
    \item \texttt{SimplexProviderFactory} — фабрика для создания объектов классов, отвечающих за построение симплекса;
    \item \texttt{HadamardMatrixIterator} — итератор для последовательного чтения матриц из файла;
    \item \texttt{LebesgueFunction} — вычисление значения функции Лебега (суммы модулей базисных многочленов) в точке;
    \item \texttt{FullEnumUpperBoundCalculator} — многопоточный перебор вершин гиперкуба;
    \item \texttt{CommandLineParser} — разбор аргументов командной строки;
    \item \texttt{Processing} — организация основных вычислительных режимов.
\end{itemize}

\subsection{Ввод и вывод данных}

Программа поддерживает два режима работы:
\begin{enumerate}
    \item Построение симплекса по матрице Адамара Сильвестра (по умолчанию);
    \item Загрузка матриц Адамара из файла (символьный или числовой формат);
\end{enumerate}

Файлы с матрицами Адамара могут быть загружены с сайта библиотеки Нила Слоуна \cite{sloane_hadamard}. Для автоматической загрузки всех доступных матриц порядка 32 разработан скрипт \texttt{download\_hadamard.sh}. Он скачивает файлы \texttt{had.32.syl}, \texttt{had.32.pal}, \texttt{had.32.t1} – \texttt{had.32.t4} и объединяет их в один файл с пустыми строками-разделителями. Кроме того, в репозитории программы в каталоге \texttt{preset} уже содержатся готовые матрицы различных размерностей, что позволяет использовать их без дополнительной загрузки.

Управление режимами работы осуществляется через аргументы командной строки:
\begin{itemize}
    \item \texttt{--dimension N} — установить размерность пространства;
    \item \texttt{--load-hadamard FILE} — загрузить матрицы Адамара из файла;
    \item \texttt{--batch DIR} — пакетная обработка всех файлов в каталоге;
    \item \texttt{--threads N} — количество потоков (0 = автоматически);
    \item \texttt{--log-interval N} — интервал логирования (количество вершин);
    \item \texttt{--range START END} — частичный перебор вершин в заданном диапазоне.
\end{itemize}


\subsection{Дополнительные проверки корректности построения проектора}

Для верификации программной реализации и достоверности полученных результатов был проведён ряд дополнительных проверок.

\subsubsection{Проверка интерполяционных свойств}

Для каждого построенного симплекса выполнялась проверка условия
\[
\lambda_j(x^{(k)}) = \delta_{jk}, \qquad j,k = 1,\dots,n+1,
\]
где $\lambda_j$ — базисные многочлены Лагранжа (барицентрические координаты), а $x^{(k)}$ — вершины симплекса. Проверка проводилась с точностью $10^{-5}$. Для всех рассмотренных матриц Адамара максимальная ошибка не превысила $10^{-10}$, что подтверждает корректность построения симплексов и вычисления базисных функций.

\subsubsection{Проверка полного перебора}

Для малых размерностей ($n = 0, 1, 3, 7, 11, 15, 19, 23, 27$) был выполнен полный перебор всех вершин гиперкуба, и полученные оценки нормы проектора были сопоставлены с известными из литературы. Результаты сравнения соответствуют приведенным в ~\ref{tab:theta_estimates}.

\subsubsection{Повторная проверка найденного максимума}

После завершения полного перебора для каждой матрицы значение нормы проектора повторно вычислялось в найденной вершине. Расхождение между значением, зафиксированным во время перебора, и значением, полученным при повторной проверке, не превышало $10^{-10}$, что подтверждает корректность синхронизации глобального максимума при многопоточном переборе.

\subsection{Вычисление нормы проектора}

Норма интерполяционного проектора вычисляется по формуле (2.3):
\[
\|P\| = \max_{x \in \operatorname{ver}(Q_n)} \sum_{j=1}^{n+1} |\lambda_j(x)|.
\]

Для нахождения максимума реализован полный перебор всех $2^{31}$ вершин гиперкуба. Каждая вершина кодируется 31-битовым целым числом, что позволяет эффективно генерировать координаты с помощью битовых операций.

Вычисления выполняются в несколько потоков (количество потоков может быть задано параметром \texttt{--threads}). Для синхронизации глобального максимума используются атомарные операции. Программа выводит прогресс выполнения каждые $10^7$ вершин с оценкой времени до завершения.

По умолчанию все результаты выводятся в стандартный поток вывода (консоль). Для сохранения результатов в файл при одновременном отображении в консоли может быть использована утилита \texttt{tee}, например:
\begin{verbatim}
 ./HadamardSimplexNorm --dimension 31 | tee results.txt
\end{verbatim}
При указании опции \texttt{--output FILE} программа самостоятельно сохраняет результаты в указанный файл.

Полный исходный код программы приведён в приложении.
% ============================================================================
\section{Результаты вычисления верхней оценки нормы}

\subsection{Параметры вычислительного эксперимента}

Вычисления проводились на следующем оборудовании и программном обеспечении:

\begin{itemize}
    \item \textbf{Процессор}: 12th Gen Intel(R) Core(TM) i5-1240P (16 ядер)
    \item \textbf{Оперативная память}: 15 ГБ
    \item \textbf{Операционная система}: Ubuntu 24.04.3 LTS
    \item \textbf{Компилятор}: GCC 13.3.0
    \item \textbf{Библиотека}: Eigen 3
    \item \textbf{Флаги компиляции}: \texttt{-O3 -march=native}
    \item \textbf{Количество потоков}: 16
    \item \textbf{Интервал логирования}: $10^7$ вершин
\end{itemize}

Полный перебор всех $2^{31}$ вершин одного гиперкуба занял 290 секунд. Средняя скорость обработки составила около 7.4 миллионов вершин в секунду.

\subsection{Результаты для размерности $n = 31$}

В результате вычислительного эксперимента для шести симплексов, построенных на основе матрицы Адамара получена следующая оценка нормы проектора:
\[
\|P\| = 5.000.
\]

\subsection{Сравнение с теоретической оценкой}

Теоретическая верхняя граница для минимальной нормы интерполяционного проектора в размерности $n = 31$ составляет:
\[
\theta_{31} \le \sqrt{32} \approx 5.656.
\]

Полученное в работе значение $\theta_{31} \le 5.000$ ниже теоретической оценки. Таким образом, наши вычисления позволили уточнить существующую оценку.
% ============================================================================

\newpage

\section*{Заключение}
\addcontentsline{toc}{section}{Заключение}

В ходе выполнения выпускной квалификационной работы была разработана программа на языке C++, предназначенная для вычисления верхних оценок минимальной нормы интерполяционного проектора $\theta_n$ для симплексов, построенных из матриц Адамара. Программа поддерживает:

\begin{itemize}
    \item построение симплекса из матрицы Адамара методом Сильвестра;
    \item загрузку матриц Адамара из файлов (символьный и числовой форматы);
    \item многопоточный полный перебор всех $2^{31}$ вершин гиперкуба;
    \item последовательная обработка нескольких файлов с данными матриц Адамара;
    \item логирование прогресса с выводом текущего результата и оценкой времени до завершения;
    \item сохранение результатов в выходной файл.
\end{itemize}

\subsection*{Полученные результаты}

В результате вычислительного эксперимента для шести симплексов, построенных по матрицам Адамара порядка $32$ различных типов, получена верхняя оценка
\[
\theta_{31} \le 5.000,
\]
Теоретическая верхняя граница для данной размерности составляет $\sqrt{32} \approx 5.656$. Таким образом, наши вычисления позволили уточнить существующую оценку.

\subsection*{Трудоёмкость анализа следующей размерности}

Следующая размерность, для которой существует матрица Адамара, — $n = 35$ (матрица порядка $36$). В этом случае гиперкуб содержит $2^{35} \approx 3.44 \times 10^{10}$ вершин. При сохранении скорости обработки $7.4 \times 10^6$ вершин в секунду полный перебор занял бы около $1.3$ часа, что выполнимо на современном компьютере. Таким образом, разработанный метод полного перебора остаётся применимым и для размерности $n = 35$.

Исходный код программы доступен в открытом репозитории по адресу:
\url{https://gitflic.ru/project/zaykov_ivan/minimal_norm_calculator.git} (или \url{https://github.com/Ivan-Zaykov/minimal_norm_calculator.git}).

% ============================================================================

\newpage

\begin{table}[h]
\centering
\caption{Количество неэквивалентных матриц Адамара}
\label{tab:hadamard_counts}
\begin{tabular}{|c|c|}
\hline
\textbf{Порядок $n$} & \textbf{Количество неэквивалентных матриц} \\
\hline
1 & 1 \\
\hline
2 & 1 \\
\hline
4 & 1 \\
\hline
8 & 1 \\
\hline
12 & 1 \\
\hline
16 & 5 \\
\hline
20 & 3 \\
\hline
24 & 60 \\
\hline
28 & 487 \\
\hline
32 & 13\,710\,027 \\
\hline
\end{tabular}
\end{table}

\newpage

\begin{table}[h]
\centering
\caption{Верхние оценки минимальной нормы интерполяционного проектора $\theta_n$ для $1 \le n \le 27$}
\label{tab:theta_estimates}
\begin{tabular}{|c|c|c|}
\hline
$n$ & $\min\|P\|$ & $\theta_n \le$ \\
\hline
1 & 1 & 1 \\
\hline
2 & 3 & $\frac{2\sqrt{5}}{5}+1 = 1.8944\ldots$ \\
\hline
3 & 2 & 2 \\
\hline
4 & $\frac{7}{3} = 2.3333\ldots$ & $\frac{3(4+\sqrt{2})}{7} = 2.3203\ldots$ \\
\hline
5 & $\frac{13}{5} = 2.6$ & $2.448804$ \\
\hline
6 & 3 & $2.60004\ldots$ \\
\hline
7 & $\frac{5}{2} = 2.5$ & $\frac{5}{2} = 2.5$ \\
\hline
8 & $\frac{22}{7} = 3.1428\ldots$ & $\frac{22}{7} = 3.1428\ldots$ \\
\hline
9 & 3 & 3 \\
\hline
10 & $\frac{19}{5} = 3.8$ & $\frac{19}{5} = 3.8\ldots$ \\
\hline
11 & 3 & 3 \\
\hline
12 & $\frac{17}{5} = 3.4$ & $\frac{17}{5} = 3.4$ \\
\hline
13 & $\frac{49}{13} = 3.7692\ldots$ & $\frac{49}{13} = 3.7692\ldots$ \\
\hline
14 & $\frac{21}{5} = 4.2$ & $\frac{21}{5} = 4.2\ldots$ \\
\hline
15 & $\frac{7}{2} = 3.5$ & $\frac{7}{2} = 3.5$ \\
\hline
16 & $\frac{21}{5} = 4.2$ & $\frac{21}{5} = 4.2$ \\
\hline
17 & $\frac{139}{34} = 4.0882\ldots$ & $\frac{139}{34} = 4.0882\ldots$ \\
\hline
18 & $\frac{95}{17} = 5.5882\ldots$ & $5.14006\ldots$ \\
\hline
19 & 4 & 4 \\
\hline
20 & $\frac{137}{29} = 4.7241\ldots$ & $4.68879\ldots$ \\
\hline
21 & $\frac{251}{50} = 5.02$ & $\frac{251}{50} = 5.02$ \\
\hline
22 & $\frac{1817}{335} = 5.4238\ldots$ & $\frac{1817}{335} = 5.4238\ldots$ \\
\hline
23 & $\frac{9}{2} = 4.5$ & $\frac{9}{2} = 4.5$ \\
\hline
24 & $\frac{103}{21} = 4.9047\ldots$ & $\frac{103}{21} = 4.9047\ldots$ \\
\hline
25 & 5 & 5 \\
\hline
26 & $\frac{474}{91} = 5.2087\ldots$ & $\frac{474}{91} = 5.2087\ldots$ \\
\hline
27 & $5.0$ & $5.0$ \\
\hline
\end{tabular}
\end{table}
\newpage

\newpage

\printbibliography

\newpage

% ============================================================================

\setupappendix
\section[Приложение \thesection]{Структура исходного кода}

Программа реализована на языке C++ (стандарт C++17). Для линейно-алгебраических операций используется библиотека Eigen. Сборка осуществляется с помощью CMake. Ниже приведена структура исходного кода:

\begin{verbatim}
├── main.cpp                     # точка входа в программу

include/
├── basis/
│   └── SimplexBasis.h           # базис Лагранжа (барицентрические координаты)
├── cli/
│   └── CommandLineParser.h      # разбор аргументов командной строки
├── hadamard_utils/
│   ├── HadamardMatrixIterator.h # итератор для чтения матриц Адамара из файла
│   └── HadamardUtils.h          # вспомогательные функции для матриц Адамара
├── processing/
│   └── Processing.h             # организация основных вычислительных режимов
├── upper_bound_calculator/
│   ├── FullEnumUpperBoundCalculator.h # многопоточный перебор вершин
│   ├── LebesgueFunction.h       # вычисление значения функции Лебега в точке
│   └── UpperBoundCalculatorInterface.h # интерфейс для вычисления верхней границы
└── vertices_provider/
    ├── HadamardFileProvider.h   # загрузка матриц Адамара из файла
    ├── ISimplexProvider.h       # интерфейс для получения вершин симплекса
    ├── SilvesterMethodProvider.h # построение симплекса методом Сильвестра
    └── SimplexProviderFactory.h # фабрика для создания провайдеров

src/
├── basis/
│   └── SimplexBasis.cpp
├── cli/
│   └── CommandLineParser.cpp
├── hadamard_utils/
│   ├── HadamardMatrixIterator.cpp
│   └── HadamardUtils.cpp
├── processing/
│   └── Processing.cpp
├── upper_bound_calculator/
│   ├── FullEnumUpperBoundCalculator.cpp
│   └── LebesgueFunction.cpp
└── vertices_provider/
    ├── HadamardFileProvider.cpp
    ├── SilvesterMethodProvider.cpp
    └── SimplexProviderFactory.cpp
\end{verbatim}

\subsection*{Компиляция и запуск}

Для компиляции программы необходимо установить библиотеку Eigen и систему сборки CMake. Команды для сборки:

\begin{verbatim}
mkdir build && cd build
cmake -DCMAKE_BUILD_TYPE=Release ..
make -j$(nproc)
\end{verbatim}

Примеры запуска:

\begin{verbatim}
# Стандартное построение (Сильвестр) для размерности 31
./HadamardSimplexNorm --dimension 31

# Пакетная обработка всех файлов в каталоге
./HadamardSimplexNorm --batch ../input --output results.txt

# Загрузка матрицы Адамара из файла
./HadamardSimplexNorm --load-hadamard input/hadamard_32.txt --dimension 31
\end{verbatim}
\end{document}